\documentclass[a4paper]{article}
\usepackage[fontset=fandol]{ctex}
\usepackage[margin=1.65cm]{geometry}
\usepackage{amsmath,amssymb}
\usepackage{graphicx}
\usepackage{hyperref}
\usepackage{longtable,booktabs}
\usepackage{xcolor}
\usepackage{enumitem}
\usepackage{listings}
\lstset{basicstyle=\ttfamily\small,breaklines=true,columns=fullflexible,keepspaces=true}
\hypersetup{colorlinks=true,linkcolor=blue,urlcolor=blue}
\setlist[itemize]{leftmargin=1.25em,itemsep=0.18em}
\setlength{\parskip}{0.45em}
\setlength{\parindent}{2em}
\sloppy
\title{第 13 讲：预训练数据：来源与数据集}
\author{CS336 Spring 2026 public videos and official course materials}
\date{}
\begin{document}
\maketitle
\tableofcontents
\newpage
\section{本章导读}
本章讨论 \textbf{预训练数据：来源与数据集}。CS336 的第一层目标不是把现成大模型当作黑箱使用，而是把 language model 从输入文本、训练目标、模型结构、数据、系统和评估这些部件重新拆开。只有拆开之后，读者才能判断一个设计选择究竟改变了什么。
设想有人说模型训练在“互联网上”。教材式回答必须拆开这句话：哪些 crawl、dump、licensed corpora、code、books、papers 和 curated datasets 被采集？谁允许使用？经过了哪些过滤？
本章对应 CS336 Spring 2026 第 13 个公开视频，并结合课程页提供的可解析资料。视频和课程页给出事实边界；正文把这些材料改写为可自学的教材章节。
阅读本章时，先理解问题为什么存在，再看定义、公式和实现形态，随后通过例题和练习检查自己是否能独立复现这条 reasoning chain。
\begin{figure}[h]
\centering
\includegraphics[width=0.58\linewidth]{source_anchor.jpg}
\caption{第 13 讲的课程图像锚点。}
\end{figure}
\subsection{本章知识路线}
\begin{itemize}
\item 本章首先说明 预训练数据：来源与数据集 为什么是 language modeling pipeline 中不可跳过的一环。
\item 其次定义本章反复使用的术语，并说明这些术语对应的计算对象或系统对象。
\item 随后用公式和伪代码表达主要机制，让概念能够落到可计算的变量上。
\item 最后给出例题、风险讨论和练习，便于读者检查自己是否真正掌握了机制。
\end{itemize}
\subsection{结构图}
下面的图用于帮助读者定位本章的核心结构。先看概念之间的依赖，再回到正文中的公式、伪代码和例题。
\begin{figure}[h]
\centering
\includegraphics[width=0.72\linewidth]{figures/source_grounded_concept_map.png}
\caption{第 13 讲概念图：本章问题、术语和机制的关系。}
\end{figure}
\begin{figure}[h]
\centering
\includegraphics[width=0.72\linewidth]{figures/source_grounded_checklist.png}
\caption{第 13 讲学习路径：从问题定义进入公式、实现、实验和练习。}
\end{figure}
\section{本章主线}
预训练数据：来源与数据集 的主线是把“训练数据”拆成来源、过滤、去重、混合和审计。模型能力不是只由参数决定，也由它反复看到的样本分布决定。
\subsection{本章要解决的核心问题}
本讲进入 data：语言模型并不是训练在“整个互联网”上，而是训练在经过 crawl、dump、extraction、filtering、dedup 和 mixture 的数据产物上。
读完这一段后，应能用一两句话说明 预训练数据：来源与数据集 要解决的瓶颈，并指出这个瓶颈发生在数据、模型、优化、系统还是评估层。
\subsection{从机制到资源账本}
数据来源有技术、法律和伦理约束。robots.txt、ToS、paywall、copyright、privacy、PII、crawler behavior 都会影响可用数据。
随后把机制写成账本：输入规模是什么，主要状态是什么，成本或质量指标怎样随变量改变。这个账本让后面的公式不再只是符号。
\subsection{从课堂讲解到可复现实验}
数据 pipeline 必须版本化。没有 provenance、filter version、dedup strategy 和 mixture weights，就无法复现实验或解释模型行为。
最后把课堂讲解转成可复现实验：固定协议，改变一个变量，记录中间量，并解释结果为何支持或反驳原来的机制判断。
\section{术语、符号与问题边界}
本章的术语横跨模型、数据和系统。读者需要先知道每个词指向的对象：有些词描述数学目标，有些词描述张量形状，有些词描述硬件瓶颈，还有一些词描述评估协议。术语边界不清，后面的公式和实现就会变成空洞的缩写。
\begin{longtable}{p{0.08\linewidth}p{0.42\linewidth}p{0.42\linewidth}}
\toprule
\textbf{\#} & \textbf{术语} & \textbf{阅读要求}\\
\midrule
1 & pretraining data（预训练数据） & 说明它如何改变训练样本或 token 分布，并给出一个会改变结论的数据边界条件。\\
2 & Common Crawl & 给出定义、一个最小例子、一个关联公式或代码位置，以及一个典型失败情形。\\
3 & data mixture（数据混合） & 说明它如何改变训练样本或 token 分布，并给出一个会改变结论的数据边界条件。\\
4 & dataset provenance（数据来源） & 说明它如何改变训练样本或 token 分布，并给出一个会改变结论的数据边界条件。\\
5 & licensing（许可证） & 给出定义、一个最小例子、一个关联公式或代码位置，以及一个典型失败情形。\\
\bottomrule
\end{longtable}
这些术语之间有依赖关系。例如 tokenization 改变 sequence length，sequence length 又改变 attention FLOPs、KV cache 和 evaluation token budget；parallelism 改变显存占用，也改变通信瓶颈和故障恢复方式。
\subsection{术语之间的关系}
在本章中，pretraining data（预训练数据） 给出问题入口，Common Crawl 描述主要机制，data mixture（数据混合） 则把机制连接到可测量的结果。读教材时应当沿着这个链条追问：输入是什么，中间状态怎样改变，输出指标为什么随之变化。
课程材料还会反复出现 data、e.g.、https、copyright、books、filtering 等词。它们不是一组同义标签，而是提示本章的不同层次：有的指对象，有的指操作，有的指约束或指标。
因此，术语表不是背诵清单，而是后文阅读的索引。遇到一个术语时，先定位它属于 representation、optimization、systems、data 还是 evaluation，再判断它改变哪一项账本。
\section{为什么数据最重要}
本节讨论 为什么数据最重要。它先限定 预训练数据：来源与数据集 要回答的基本问题：哪些现象需要解释，哪些变量可以测量，哪些结论必须回到 样本数量、去重率、过滤阈值、数据混合比例和 contamination 风险 中检验。
\subsection{课程目标与问题边界}
视频开场说：给定数据后大家已知道如何训练，接下来问题是应该训练在什么数据上。
数据决定模型看到的世界、语言、风格、知识和偏见，也决定 downstream 能力的上限。
课程把互联网数据看作从 live services 到 dump/crawl 再到 processing pipeline 的产物。
为什么数据最重要 把 预训练数据：来源与数据集 连接到训练分布。接下来要看的不是语料名称，而是样本来源、处理顺序和版本记录怎样改变模型可学习的证据。
\subsection{为什么数据最重要 的机制}
视频开场说：给定数据后大家已知道如何训练，接下来问题是应该训练在什么数据上。 在 为什么数据最重要 中，核心对象是样本、来源和分布状态：数据管线每改变一次样本集合，模型实际学习的分布也随之改变。
数据决定模型看到的世界、语言、风格、知识和偏见，也决定 downstream 能力的上限。 因而同一个方法在不同规模下可能改变不同的变量，尤其是样本数量、去重率、过滤阈值、数据混合比例和 contamination 风险。比较方法时，必须同时给出问题规模和实验条件。
课程把互联网数据看作从 live services 到 dump/crawl 再到 processing pipeline 的产物。 可检验性来自 provenance 记录、过滤前后计数、抽样审计和 held-out contamination 检查；没有这些记录，数据结论无法复现。
因此，为什么数据最重要 的学习目标是建立一条数据谱系：读者应能从原始来源追到训练 token，并说明每个处理步骤改变了什么分布。
\subsection{资源、效率与效果的关系}
\[\mathcal{D}_{\text{train}}\sim \sum_{k=1}^{K} w_k\,\mathcal{D}_k\]
符号说明：\(\mathcal{D}_k\) 是第 \(k\) 个数据源，\(w_k\) 是混合权重；预训练分布是人为设计的混合。
这个关系式刻画 为什么数据最重要 对数据分布的影响。读者应关心哪些变量来自采集过程，哪些变量来自过滤或混合策略，哪些变量只能通过抽样审计估计。
在数据实验中，去重率、过滤通过率、语言比例和抽样人工判读都是 proxy。它们不能单独代表数据质量，必须和来源记录、版本号以及污染检查一起解释。
\subsection{把课程目标写成检查流程}
把 为什么数据最重要 写成程序时，要明确每一步怎样改写样本集合。下面的伪代码强调输入、过滤、统计和输出之间的关系。
\begin{lstlisting}[language=Python]
sources = [web, books, code, wikipedia, papers]
for source in sources:
    collect_with_provenance(source)
    normalize_and_filter(source)
\end{lstlisting}
阅读这段流程时，重点看 为什么数据最重要 怎样改变样本集合。每个过滤、去重或混合步骤都应留下计数和抽样检查，使 样本数量、去重率、过滤阈值、数据混合比例和 contamination 风险 能够被复核。
\subsection{边界条件与常见误解}
\begin{itemize}
\item 数据不是从天上掉下来；每个来源都有采集、许可和偏差问题。
\item 更多 token 不等于更好，低质量 token 会稀释训练预算。
\item 小样本抽样只能发现部分数据问题，不能证明语料整体无污染或无偏差。
\item 数据版本、许可状态和处理脚本如果缺失，后续 evaluation 与 scaling 结论都会失去可复现基础。
\end{itemize}
\subsection{与本章其他概念的关系}
为什么数据最重要 连接本章的数据来源和模型行为：定义说明样本如何进入语料，公式刻画分布变化，伪代码展示处理顺序，例题检验数据决策怎样影响训练结论。
\subsection{怎样检验 这一课程目标}
检验 这一课程目标 时，先把抽象说法落到样本、来源和分布状态上。与 data、Common Crawl、source、dataset、provenance 相关的对象必须能被构造、打印、计数或评分；否则实验只能得到描述性观察，不能得到可复现结论。
一个合适的小实验应当保留本节机制，同时让 样本数量、去重率、过滤阈值、数据混合比例和 contamination 风险 中至少一个量发生可解释变化。输入规模不必逼真；关键是读者能手算或打印关键中间量，并说明哪个变量触发了结果变化。
随后把公式说明转成可记录的数量关系。符号说明：\textbackslash{}(\textbackslash{}mathcal\{D\}\_k\textbackslash{}) 是第 \textbackslash{}(k\textbackslash{}) 个数据源，\textbackslash{}(w\_k\textbackslash{}) 是混合权重；预训练分布是人为设计的混合。 如果数量只能间接观测，就必须说明使用了什么 proxy；如果使用 benchmark 或 reward，也要说明协议和随机性。
数据消融通常一次只改变过滤阈值、去重策略、混合权重或数据版本。若同时更换 tokenizer、语料和评估集，就无法知道改进来自哪一步。
最后检查失败条件。工程判断往往在反例中变清楚：数据不是从天上掉下来；每个来源都有采集、许可和偏差问题；更多 token 不等于更好，低质量 token 会稀释训练预算。复习时可以把这些限制改写成反例测试，观察它们怎样改变 样本数量、去重率、过滤阈值、数据混合比例和 contamination 风险。
把实验结论放回 预训练数据：来源与数据集 时，要说明 为什么数据最重要 改变的是样本分布、数据质量、污染风险还是可复现性。能追到这些来源，才说明读者真正理解数据机制。
\subsection{概念辨析：为什么数据最重要}
data、Common Crawl、source、dataset、provenance 常一起出现在数据管线中，但它们分别指来源、处理动作、版本记录或风险标记。读者应先判断它们改变的是样本集合还是审计证据。
其中，data（数据）是模型能力、偏差和法律风险的主要来源；data（数据）是模型能力、偏差和法律风险的主要来源。
区分这些词时，可以先问它们是否改变 样本数量、去重率、过滤阈值、数据混合比例和 contamination 风险。一个术语如果不能对应到变量、状态、代码路径或评估协议，就还没有形成可操作含义。
围绕 为什么数据最重要 复习时，可以构造一小批文档：保留来源和版本，运行一个处理步骤，再比较样本计数和抽样质量。
\subsection{为什么数据最重要 的小结}
\begin{itemize}
\item 讨论 为什么数据最重要 时，应同时报告问题规模、实验设置和主要观测变量，尤其是 样本数量、去重率、过滤阈值、数据混合比例和 contamination 风险。
\item 如果方法声称数据更好，要说明过滤、去重、混合或合成步骤改变了哪一部分分布，并检查 benchmark contamination。
\item 数据结论需要抽样审计和版本记录；没有 provenance 的语料变化不能支持可引用结论。
\end{itemize}
\section{Web、code、books 与 curated datasets}
现在进入 Web、code、books 与 curated datasets。这一节关心训练样本如何形成、被筛选并进入混合分布；它直接决定模型看到什么，也决定 evaluation 是否可能被污染。
\subsection{为什么要讨论 Web、code、books 与 curated datasets}
Common Crawl 提供大规模网页原料，但需要 HTML extraction、language ID、质量过滤和去重。
code 数据对编程能力关键，但有许可证、重复、生成数据污染和安全风险。
books、Wikipedia、论文和问答数据更 curated，但规模、许可和代表性不同。
数据部分的核心问题是分布控制。采集、过滤、去重和混合并不会只改变文件大小，它们会改变模型看到的语言、知识、偏差和 benchmark contamination 风险。
\subsection{Web、code、books 与 curated datasets 的机制}
Common Crawl 提供大规模网页原料，但需要 HTML extraction、language ID、质量过滤和去重。 在 Web、code、books 与 curated datasets 中，核心对象是样本、来源和分布状态：数据管线每改变一次样本集合，模型实际学习的分布也随之改变。
code 数据对编程能力关键，但有许可证、重复、生成数据污染和安全风险。 因而同一个方法在不同规模下可能改变不同的变量，尤其是样本数量、去重率、过滤阈值、数据混合比例和 contamination 风险。比较方法时，必须同时给出问题规模和实验条件。
books、Wikipedia、论文和问答数据更 curated，但规模、许可和代表性不同。 可检验性来自 provenance 记录、过滤前后计数、抽样审计和 held-out contamination 检查；没有这些记录，数据结论无法复现。
因此，Web、code、books 与 curated datasets 的学习目标是建立一条数据谱系：读者应能从原始来源追到训练 token，并说明每个处理步骤改变了什么分布。
\subsection{Web、code、books 与 curated datasets 的数据分布视角}
\[\text{value/source}=\frac{\Delta \text{eval score}}{\text{tokens}\times\text{processing cost}}\]
符号说明：这个比值是延伸解释，用来表达单位 token 和单位处理成本的边际价值，不是官方唯一指标。
这个关系式刻画 Web、code、books 与 curated datasets 对数据分布的影响。读者应关心哪些变量来自采集过程，哪些变量来自过滤或混合策略，哪些变量只能通过抽样审计估计。
在数据实验中，去重率、过滤通过率、语言比例和抽样人工判读都是 proxy。它们不能单独代表数据质量，必须和来源记录、版本号以及污染检查一起解释。
\subsection{Web、code、books 与 curated datasets 在数据管线中的实现}
把 Web、code、books 与 curated datasets 写成程序时，要明确每一步怎样改写样本集合。下面的伪代码强调输入、过滤、统计和输出之间的关系。
\begin{lstlisting}[language=Python]
for document in crawl:
    text = extract_main_content(document.html)
    if language_ok(text) and quality_ok(text):
        write_jsonl(text, metadata=document.metadata)
\end{lstlisting}
阅读这段流程时，重点看 Web、code、books 与 curated datasets 怎样改变样本集合。每个过滤、去重或混合步骤都应留下计数和抽样检查，使 样本数量、去重率、过滤阈值、数据混合比例和 contamination 风险 能够被复核。
\subsection{边界条件与常见误解}
\begin{itemize}
\item 网页抽取错误会把导航栏、广告、脚本或重复模板带入训练。
\item 不同来源的 license 和 consent 不能在技术 pipeline 中被忽略。
\item 小样本抽样只能发现部分数据问题，不能证明语料整体无污染或无偏差。
\item 数据版本、许可状态和处理脚本如果缺失，后续 evaluation 与 scaling 结论都会失去可复现基础。
\end{itemize}
\subsection{与本章其他概念的关系}
Web、code、books 与 curated datasets 连接本章的数据来源和模型行为：定义说明样本如何进入语料，公式刻画分布变化，伪代码展示处理顺序，例题检验数据决策怎样影响训练结论。
\subsection{怎样检验 Web、code、books 与 curated datasets}
检验 Web、code、books 与 curated datasets 时，先把抽象说法落到样本、来源和分布状态上。与 web、code、books、Wikipedia、license 相关的对象必须能被构造、打印、计数或评分；否则实验只能得到描述性观察，不能得到可复现结论。
一个合适的小实验应当保留本节机制，同时让 样本数量、去重率、过滤阈值、数据混合比例和 contamination 风险 中至少一个量发生可解释变化。输入规模不必逼真；关键是读者能手算或打印关键中间量，并说明哪个变量触发了结果变化。
随后把公式说明转成可记录的数量关系。符号说明：这个比值是延伸解释，用来表达单位 token 和单位处理成本的边际价值，不是官方唯一指标。 如果数量只能间接观测，就必须说明使用了什么 proxy；如果使用 benchmark 或 reward，也要说明协议和随机性。
数据消融通常一次只改变过滤阈值、去重策略、混合权重或数据版本。若同时更换 tokenizer、语料和评估集，就无法知道改进来自哪一步。
最后检查失败条件。工程判断往往在反例中变清楚：网页抽取错误会把导航栏、广告、脚本或重复模板带入训练；不同来源的 license 和 consent 不能在技术 pipeline 中被忽略。复习时可以把这些限制改写成反例测试，观察它们怎样改变 样本数量、去重率、过滤阈值、数据混合比例和 contamination 风险。
把实验结论放回 预训练数据：来源与数据集 时，要说明 Web、code、books 与 curated datasets 改变的是样本分布、数据质量、污染风险还是可复现性。能追到这些来源，才说明读者真正理解数据机制。
\subsection{概念辨析：Web、code、books 与 curated datasets}
web、code、books、Wikipedia、license 常一起出现在数据管线中，但它们分别指来源、处理动作、版本记录或风险标记。读者应先判断它们改变的是样本集合还是审计证据。
其中，decode（解码）逐 token 生成，常受 memory bandwidth 和 KV cache 影响。
区分这些词时，可以先问它们是否改变 样本数量、去重率、过滤阈值、数据混合比例和 contamination 风险。一个术语如果不能对应到变量、状态、代码路径或评估协议，就还没有形成可操作含义。
围绕 Web、code、books 与 curated datasets 复习时，可以构造一小批文档：保留来源和版本，运行一个处理步骤，再比较样本计数和抽样质量。
\subsection{Web、code、books 与 curated datasets 的小结}
\begin{itemize}
\item 讨论 Web、code、books 与 curated datasets 时，应同时报告问题规模、实验设置和主要观测变量，尤其是 样本数量、去重率、过滤阈值、数据混合比例和 contamination 风险。
\item 如果方法声称数据更好，要说明过滤、去重、混合或合成步骤改变了哪一部分分布，并检查 benchmark contamination。
\item 数据结论需要抽样审计和版本记录；没有 provenance 的语料变化不能支持可引用结论。
\end{itemize}
\section{数据集记录与可复现性}
现在进入 数据集记录与可复现性。这一节关心训练样本如何形成、被筛选并进入混合分布；它直接决定模型看到什么，也决定 evaluation 是否可能被污染。
\subsection{为什么要讨论 数据集记录与可复现性}
课程强调开放模型与开放数据对信任和创新的作用；数据集 provenance 是可复现训练的基础。
每个数据版本应记录来源 URL、抓取时间、过滤器版本、去重策略、tokenizer 和混合权重。
数据集本身应当像模型 checkpoint 一样版本化，否则 evaluation 和 scaling 结果难以复现。
数据集记录与可复现性 把 预训练数据：来源与数据集 连接到训练分布。接下来要看的不是语料名称，而是样本来源、处理顺序和版本记录怎样改变模型可学习的证据。
\subsection{数据集记录与可复现性 的机制}
课程强调开放模型与开放数据对信任和创新的作用；数据集 provenance 是可复现训练的基础。 在 数据集记录与可复现性 中，核心对象是样本、来源和分布状态：数据管线每改变一次样本集合，模型实际学习的分布也随之改变。
每个数据版本应记录来源 URL、抓取时间、过滤器版本、去重策略、tokenizer 和混合权重。 因而同一个方法在不同规模下可能改变不同的变量，尤其是样本数量、去重率、过滤阈值、数据混合比例和 contamination 风险。比较方法时，必须同时给出问题规模和实验条件。
数据集本身应当像模型 checkpoint 一样版本化，否则 evaluation 和 scaling 结果难以复现。 可检验性来自 provenance 记录、过滤前后计数、抽样审计和 held-out contamination 检查；没有这些记录，数据结论无法复现。
因此，数据集记录与可复现性 的学习目标是建立一条数据谱系：读者应能从原始来源追到训练 token，并说明每个处理步骤改变了什么分布。
\subsection{数据集记录与可复现性 的数学表达}
\[\text{dataset version}=f(\text{sources},\text{filters},\text{dedup},\text{tokenizer},\text{mixture})\]
符号说明：版本不是文件名，而是由完整 pipeline 决定的可审计对象。
这个关系式刻画 数据集记录与可复现性 对数据分布的影响。读者应关心哪些变量来自采集过程，哪些变量来自过滤或混合策略，哪些变量只能通过抽样审计估计。
在数据实验中，去重率、过滤通过率、语言比例和抽样人工判读都是 proxy。它们不能单独代表数据质量，必须和来源记录、版本号以及污染检查一起解释。
\subsection{从 数据集记录与可复现性 到程序}
把 数据集记录与可复现性 写成程序时，要明确每一步怎样改写样本集合。下面的伪代码强调输入、过滤、统计和输出之间的关系。
\begin{lstlisting}[language=Python]
manifest = {
    'source': url,
    'crawl_time': timestamp,
    'filters': filter_hash,
    'tokenizer': tokenizer_hash,
}
\end{lstlisting}
阅读这段流程时，重点看 数据集记录与可复现性 怎样改变样本集合。每个过滤、去重或混合步骤都应留下计数和抽样检查，使 样本数量、去重率、过滤阈值、数据混合比例和 contamination 风险 能够被复核。
\subsection{边界条件与常见误解}
\begin{itemize}
\item 只发布 token count 不足以复现数据。
\item 数据删除请求、法律要求和安全过滤会让数据版本随时间变化。
\item 小样本抽样只能发现部分数据问题，不能证明语料整体无污染或无偏差。
\item 数据版本、许可状态和处理脚本如果缺失，后续 evaluation 与 scaling 结论都会失去可复现基础。
\end{itemize}
\subsection{与本章其他概念的关系}
数据集记录与可复现性 连接本章的数据来源和模型行为：定义说明样本如何进入语料，公式刻画分布变化，伪代码展示处理顺序，例题检验数据决策怎样影响训练结论。
\subsection{怎样检验 数据集记录与可复现性}
检验 数据集记录与可复现性 时，先把抽象说法落到样本、来源和分布状态上。与 manifest、version、provenance、filter、tokenizer 相关的对象必须能被构造、打印、计数或评分；否则实验只能得到描述性观察，不能得到可复现结论。
一个合适的小实验应当保留本节机制，同时让 样本数量、去重率、过滤阈值、数据混合比例和 contamination 风险 中至少一个量发生可解释变化。输入规模不必逼真；关键是读者能手算或打印关键中间量，并说明哪个变量触发了结果变化。
随后把公式说明转成可记录的数量关系。符号说明：版本不是文件名，而是由完整 pipeline 决定的可审计对象。 如果数量只能间接观测，就必须说明使用了什么 proxy；如果使用 benchmark 或 reward，也要说明协议和随机性。
数据消融通常一次只改变过滤阈值、去重策略、混合权重或数据版本。若同时更换 tokenizer、语料和评估集，就无法知道改进来自哪一步。
最后检查失败条件。工程判断往往在反例中变清楚：只发布 token count 不足以复现数据；数据删除请求、法律要求和安全过滤会让数据版本随时间变化。复习时可以把这些限制改写成反例测试，观察它们怎样改变 样本数量、去重率、过滤阈值、数据混合比例和 contamination 风险。
把实验结论放回 预训练数据：来源与数据集 时，要说明 数据集记录与可复现性 改变的是样本分布、数据质量、污染风险还是可复现性。能追到这些来源，才说明读者真正理解数据机制。
\subsection{概念辨析：数据集记录与可复现性}
manifest、version、provenance、filter、tokenizer 常一起出现在数据管线中，但它们分别指来源、处理动作、版本记录或风险标记。读者应先判断它们改变的是样本集合还是审计证据。
区分这些词时，可以先问它们是否改变 样本数量、去重率、过滤阈值、数据混合比例和 contamination 风险。一个术语如果不能对应到变量、状态、代码路径或评估协议，就还没有形成可操作含义。
围绕 数据集记录与可复现性 复习时，可以构造一小批文档：保留来源和版本，运行一个处理步骤，再比较样本计数和抽样质量。
\subsection{数据集记录与可复现性 的小结}
\begin{itemize}
\item 讨论 数据集记录与可复现性 时，应同时报告问题规模、实验设置和主要观测变量，尤其是 样本数量、去重率、过滤阈值、数据混合比例和 contamination 风险。
\item 如果方法声称数据更好，要说明过滤、去重、混合或合成步骤改变了哪一部分分布，并检查 benchmark contamination。
\item 数据结论需要抽样审计和版本记录；没有 provenance 的语料变化不能支持可引用结论。
\end{itemize}
\section{综合讨论：从局部机制到完整系统}
本章的主题是 预训练数据：来源与数据集，但它不应被读成几个相邻知识点的拼接。为什么数据最重要、Web、code、books 与 curated datasets、数据集记录与可复现性 共同回答同一个问题：语言模型系统中哪些对象可以被定义，哪些成本可以被估算，哪些结论必须通过实验或评估来确认。
这些对象包括 pretraining data（预训练数据）、Common Crawl、data mixture（数据混合）、dataset provenance（数据来源）、licensing（许可证）。读者可以把它们看成一组接口：有的接口连接文本和 token，有的连接模型结构和张量计算，有的连接数据样本和训练目标，有的连接离线评估和在线服务。接口一旦改变，后面的成本、错误类型和优化空间也会随之改变。
\subsection{为什么数据最重要 在系统中的位置}
为什么数据最重要 首先提供一个局部解释：视频开场说：给定数据后大家已知道如何训练，接下来问题是应该训练在什么数据上。 这类解释的价值在于把直觉压缩成可操作对象。一旦对象明确，读者就可以追问它的输入、输出、中间状态和失败条件，而不是只记住方法名。
其次，为什么数据最重要 会改变至少一个资源或评估变量。它可能改变 sequence length、activation size、memory bandwidth、communication volume、data mixture、reward distribution 或 benchmark protocol。这些变量在小例子里可以手算，在真实训练和服务系统里则需要 profiling、logging 和 held-out evaluation。
最后，为什么数据最重要 不能脱离边界条件理解。一个关键 caveat 是：数据不是从天上掉下来；每个来源都有采集、许可和偏差问题。。把 caveat 写出来不是为了削弱结论，而是为了说明结论在哪些输入、硬件、数据和评估条件下才成立。
\subsection{Web、code、books 与 curated datasets 在系统中的位置}
Web、code、books 与 curated datasets 首先提供一个局部解释：Common Crawl 提供大规模网页原料，但需要 HTML extraction、language ID、质量过滤和去重。 这类解释的价值在于把直觉压缩成可操作对象。一旦对象明确，读者就可以追问它的输入、输出、中间状态和失败条件，而不是只记住方法名。
其次，Web、code、books 与 curated datasets 会改变至少一个资源或评估变量。它可能改变 sequence length、activation size、memory bandwidth、communication volume、data mixture、reward distribution 或 benchmark protocol。这些变量在小例子里可以手算，在真实训练和服务系统里则需要 profiling、logging 和 held-out evaluation。
最后，Web、code、books 与 curated datasets 不能脱离边界条件理解。一个关键 caveat 是：网页抽取错误会把导航栏、广告、脚本或重复模板带入训练。。把 caveat 写出来不是为了削弱结论，而是为了说明结论在哪些输入、硬件、数据和评估条件下才成立。
\subsection{数据集记录与可复现性 在系统中的位置}
数据集记录与可复现性 首先提供一个局部解释：课程强调开放模型与开放数据对信任和创新的作用；数据集 provenance 是可复现训练的基础。 这类解释的价值在于把直觉压缩成可操作对象。一旦对象明确，读者就可以追问它的输入、输出、中间状态和失败条件，而不是只记住方法名。
其次，数据集记录与可复现性 会改变至少一个资源或评估变量。它可能改变 sequence length、activation size、memory bandwidth、communication volume、data mixture、reward distribution 或 benchmark protocol。这些变量在小例子里可以手算，在真实训练和服务系统里则需要 profiling、logging 和 held-out evaluation。
最后，数据集记录与可复现性 不能脱离边界条件理解。一个关键 caveat 是：只发布 token count 不足以复现数据。。把 caveat 写出来不是为了削弱结论，而是为了说明结论在哪些输入、硬件、数据和评估条件下才成立。
\subsection{把本章知识用于读论文和复现实验}
读相关论文时，可以把本章变成一个检查表。第一，论文是否清楚定义了输入规模、模型规模、数据来源和评估协议；第二，论文声称的改进落在哪个变量上；第三，作者是否报告了会让结论失效的设置；第四，实验是否能分离模型结构、数据质量、优化设置和系统实现的影响。
复现实验时，也应避免直接追求大规模。先用最小程序复现一个公式、一个伪代码分支或一个 profile 现象，再扩大输入规模。这样做的原因很简单：如果小规模程序无法解释中间量，大规模训练只会把错误隐藏在日志和随机波动里。
把这些检查合在一起，本章给出的不是一个固定答案，而是一种阅读 CS336 的方法：每个模型机制都要落到变量，每个变量都要落到测量，每个测量都要说明协议，每个协议都要承认边界。
因此，本章中的任何一句结论都可以被改写成一个可引用的完整句式：在给定数据、模型、硬件和评估协议下，某个机制改变了某个变量，并通过某个观测指标表现出来。这样的句子比单独背诵术语更长，但它包含了自学、复现和引用时真正需要的信息。
把本章用于自学时，可以按三遍阅读。第一遍只追踪对象和定义，确认每个术语到底指向文本、token、张量、样本、请求还是评估记录。第二遍追踪公式和伪代码，确认每一步改变了哪些状态。第三遍追踪实验和 caveat，确认哪些结论只在特定规模、数据或硬件条件下成立。
把本章用于复习时，则应反过来操作：先从一个失败案例或异常指标出发，再回到相应机制。loss 曲线异常可能指向数据、优化或模型容量；吞吐异常可能指向 kernel、通信或调度；benchmark 异常可能指向 contamination、prompt protocol 或采样设置。能够从异常反推机制，是掌握本章的实用标准。
\section{例题：把概念变成可计算检查}
\subsection{从网页到训练样本的 provenance}
一个网页进入训练集，要经过 crawl、HTML extraction、language ID、quality filter 和 dedup。
\begin{itemize}
\item 每一步都应写入数据清单（manifest），而不是只留下最终文本。
\item 过滤器版本和阈值会改变数据分布。
\item 许可、robots、PII 和删除请求会影响能否持续使用。
\end{itemize}
\textbf{结论。} 数据集是 pipeline 的产物，不是文件夹里若干 `.jsonl` 的集合。
这个例题的目的不是替代完整工程实现，而是给出一种最小推理形式：把问题转成变量，写出近似公式或伪代码，再说明近似何时失效。
\section{实验设计：从小例子到可复现结论}
学习 预训练数据：来源与数据集 不能停在概念层。一个教材化结论至少需要四件事：可构造的输入、可观测的中间量、可比较的指标，以及清楚的失败条件。下面的实验设计不是要求训练大模型，而是要求读者用小规模程序复现本章机制。
\subsection{最小输入}
先构造只触发本章核心机制的小输入。例如 tokenizer 章节可以用几个包含 rare words 和 Unicode 字符的字符串；GPU 章节可以用一个 elementwise kernel 和一个 GEMM；data 章节可以用十几篇近重复文档。小输入的价值在于它让错误可见。
\subsection{资源账本}
对同一输入写下 FLOPs、bytes moved、显存峰值、通信量或样本数量的估算。估算不需要一开始就精确，但必须说明单位和近似假设。随后用 profiler、benchmark、validation loss 或人工检查去验证估算是否同量级。
\subsection{单变量消融}
一次只改变一个因素，例如 vocabulary size、head 数、batch size、filter threshold、reward scale、decoding temperature 或 GPU 数量。若多个变量同时变化，实验只能说明系统变了，不能说明机制是什么。
\subsection{失败条件}
最后要刻意触发失败：极端 context length、错误 dtype、过小 batch、污染 benchmark、低质量数据或不可靠 verifier。失败条件常常比成功曲线更能说明一个方法的边界。
\subsection{记录与报告}
实验记录应能让另一个读者复现同一结论。最少要写清楚输入构造、代码版本、随机种子、硬件、batch shape、数据版本、评价指标和异常处理方式。报告结论时不要只写“更快”或“更好”，而要说明快在哪里、好在哪个指标上，以及是否牺牲了内存、稳定性、数据质量或可解释性。
常见报告错误是把局部现象写成普遍规律。例如一次小模型实验里的 loss 改善，不等于更大模型或不同数据分布上也会改善；一次 kernel benchmark 的吞吐提升，不等于端到端 serving latency 一定降低。教材中的实验报告应始终把结论限定在可复现条件内。
\subsection{从小实验到大结论}
小实验的作用不是替代课程中的完整训练或系统实现，而是建立一条可检查的推理链。先在小输入上确认变量方向，再在中等规模上确认数量级，最后才讨论是否能外推到更大模型、更多 token 或更复杂 workload。缺少中间层级时，大结论通常只是在复述直觉。
因此，实验报告最好同时包含正例和反例：正例说明机制在受控条件下怎样工作，反例说明条件稍变时哪里失效。对 CS336 这样的课程，反例尤其重要，因为 language model 的错误常来自多个层面叠加：数据分布、优化稳定性、硬件瓶颈、评估协议和服务负载都可能同时改变观察结果。
还要记录资料版本与复现边界。公开视频、课程页、配套代码和 PDF 可能在不同时间更新；复现实验应写明使用的材料版本，并说明哪些结论来自课程事实，哪些是为了帮助理解而加入的解释性推导。这样才能把学习笔记变成可检查的技术记录。
若复现实验与课程结论不一致，首先检查输入、版本和评估协议，而不是立即修改模型假设。
\subsection{把本章重新读成一条证据链}
完成小实验之后，可以把 预训练数据：来源与数据集 重新读成一条证据链：先用“为什么数据最重要”确定问题入口，再用“Web、code、books 与 curated datasets”解释主要机制，最后用“数据集记录与可复现性”检查边界或外推条件。这种读法比按页背诵更接近教材训练，因为它要求每个结论都能说明前提、变量和证据。
本章反复使用的术语包括 pretraining data（预训练数据）、Common Crawl、data mixture（数据混合）、dataset provenance（数据来源）、licensing（许可证）。复习时不要只写定义，而要为每个术语补上一句操作性说明：它在课程代码、公式、数据记录或评估协议中对应什么对象；如果这个对象被删除、替换或缩放，哪一个观测量会变化。
最后，用一个反例结束复习。反例可以是资源估算不准、数据污染、reward 失真、benchmark 解析失败、长上下文显存爆炸或 scaling 外推失效。能解释反例的章节，才真正具备自学和引用价值；不能解释反例的章节，只是把课程材料换了一种排版。
引用本章结论时，还应保留来源边界：哪些说法直接来自视频、课程页、slides 或代码，哪些是为了串联教材叙述而加入的解释性推导。这个边界不会削弱结论，反而让读者知道何处可以复核，何处需要在自己的实验中重新验证。如果后续课程材料更新，也可以沿着同一证据链替换来源，而不必重写整章结构。
\section{课程资料与回看路径}
教材正文已经把主要概念、公式和实现逻辑组织成连续叙述；本节只提供回到课程材料的阅读路径。读者复习时可以先完成前文练习，再按下面的时间段或资料组核对细节。
\subsection{视频回看路径}
\begin{longtable}{p{0.20\linewidth}p{0.46\linewidth}p{0.26\linewidth}}
\toprule
\textbf{时间} & \textbf{关键词} & \textbf{复习任务}\\
\midrule
00:00:04-00:06:20 & data, training, there's, data., that's, talk, train & 回看定义、例子、推导或 caveat 在该段中的位置。\\
00:19:25-00:26:05 & fair, copyright, license, creative, it's, that's, don't & 回看定义、例子、推导或 caveat 在该段中的位置。\\
00:40:03-00:46:20 & data, it's, github, books, there's, crawl, code & 回看定义、例子、推导或 caveat 在该段中的位置。\\
00:53:24-01:00:55 & data, think, books3, wikipedia, books, trained, quality & 回看定义、例子、推导或 caveat 在该段中的位置。\\
01:14:57-01:21:49 & data, it's, permissively, think, there's, very, whether & 回看定义、例子、推导或 caveat 在该段中的位置。\\
\bottomrule
\end{longtable}
这些时间段用于定位视频中的讲解顺序。严肃引用时，应回到视频片段确认讲者表述、板书或幻灯片内容。
回看视频时，不必逐字抄录字幕。更有效的方法是把每个时间段判定为定义、例子、推导、代码解释或 caveat，再把它放回本章对应小节。
\subsection{课程资料阅读路径}
\begin{longtable}{p{0.14\linewidth}p{0.48\linewidth}p{0.30\linewidth}}
\toprule
\textbf{资料组} & \textbf{标题/页块} & \textbf{关键词}\\
\midrule
official\_01 & Lecture 13: Data I & data, train, training, live, language, pre-training, mid-training\\
official\_03 & - Threshold for copyright is extremely low (e.g., your website is copyrighted) & copyright, work, fair, license, training, copyrighted, over\\
official\_05 & - Repository: download through git protocol (rather than scraping the GitHub website) & https, github, metadata, arxiv, books, language, crawl\\
official\_08 & - ... & data, tokens, trained, classifier, result, remove, permissively\\
official\_10 & Reproduced by Together's RedPajama v1 & function, refinedweb, dolma, stack, reproduced, together, redpajama\\
\bottomrule
\end{longtable}
课程资料通常给出更稳定的标题、公式、图或代码结构；视频则补充动机和口头解释。两者合起来构成本章的事实边界。
阅读资料时，可以把每个资料组拆成三个对象：一个定义，一个公式或伪代码动作，一个边界条件。这样比逐字翻译页面或脚本更容易形成可用知识。
\section{实践说明、历史位置与风险}
预训练数据：来源与数据集 在 CS336 的课程结构中连接基础机制和规模化问题。基础机制解释方法为什么成立；规模化问题检验它在更大模型、更长上下文、更复杂数据或真实 serving workload 下是否仍成立。
历史位置不等于模型年表。更重要的是识别问题如何迁移：早期方法解决小规模建模问题，现代方法常常解决同一问题在大规模数据、GPU 集群、长上下文或人类偏好约束下的新形态。
\begin{itemize}
\item 资料版本限制（version risk）：课程页、公开视频和配套资料可能在学期中更新；引用时应注明访问日期或资料版本。
\item 测量口径限制（measurement risk）：FLOPs、tokens/s、benchmark score、reward 等数字只有在协议完整时可比较。
\item 规模外推限制（scaling risk）：小模型或小 batch 的机制不必然外推到 frontier-scale。
\item 数据分布限制（data risk）：数据来源、过滤和去重策略会改变模型行为，也会改变 evaluation contamination 风险。
\end{itemize}
\section{复习要点}
下面的问题把本章内容压缩为可反复检查的条目。每个条目都应能回到前文的解释、公式、伪代码或例题。
\begin{longtable}{p{0.16\linewidth}p{0.38\linewidth}p{0.38\linewidth}}
\toprule
\textbf{条目} & \textbf{对象} & \textbf{检查动作}\\
\midrule
条目 1 & pretraining data（预训练数据） & 定义它；写出一个最小例子；说明一个 failure mode；指出它影响哪个资源或评估账本。\\
条目 2 & Common Crawl & 定义它；写出一个最小例子；说明一个 failure mode；指出它影响哪个资源或评估账本。\\
条目 3 & data mixture（数据混合） & 定义它；写出一个最小例子；说明一个 failure mode；指出它影响哪个资源或评估账本。\\
条目 4 & dataset provenance（数据来源） & 定义它；写出一个最小例子；说明一个 failure mode；指出它影响哪个资源或评估账本。\\
条目 5 & licensing（许可证） & 定义它；写出一个最小例子；说明一个 failure mode；指出它影响哪个资源或评估账本。\\
\bottomrule
\end{longtable}
\subsection{综合自测}
\begin{itemize}
\item 我能否不用课程原句，独立解释本章中心问题？
\item 我能否把本章至少一个公式改写成可运行的估算函数？
\item 我能否指出一个看似改进但实际转移成本的方法？
\item 我能否回到视频和课程资料，找到支持本章关键论断的位置？
\item 我能否设计一个最小 ablation，验证本章一个 caveat？
\end{itemize}
\section{总结与延伸}
\subsection{本章小结}
\begin{itemize}
\item 预训练数据：来源与数据集 的核心不是记忆术语，而是理解对象、变量和机制之间的关系。
\item 视频给出讲解顺序，课程资料提供可核对的公式、图、代码或页面结构。
\item 复习时应把每个术语连接到至少一个变量、一个实现动作和一个失败模式。
\end{itemize}
\subsection{延伸解释}
本章中的延伸解释用于连接课程材料中的概念，例如说明公式里的 proxy、解释系统 tradeoff，或把多个材料片段串成一个完整推理链。若需要严肃引用，应回到课程视频和配套资料核对原始表述。
\subsection{练习}
\begin{itemize}
\item 定义题：用中英双语解释本章任意五个重要术语，并说明它们属于 compute、memory、data、optimization、evaluation 还是 deployment 账本。
\item 推导题：选择本章一个公式，逐个解释符号、单位和近似假设，并说明哪个变量最难在真实系统中测量。
\item 实现题：把本章伪代码改写成一个最小 Python sanity check，不要求训练大模型，但要输出可检查的中间量。
\item 资料题：从“课程资料与回看路径”中选择一个视频时间段，回到原始视频核对讲者例子，并记录是否存在字幕误差。
\item 评估题：为本章方法设计一个 ablation 或 benchmark，说明控制变量、数据版本、指标和 failure criteria。
\item 边界题：说明本章一个结论在更长 context、更大 batch、更多 GPU 或更复杂数据分布下可能如何失效。
\item 综合题：把本章内容和前一章或后一章连接起来，写出一个端到端训练/推理 pipeline 中的依赖关系。
\item 批判题：指出一个容易被营销材料夸大的说法，并用本章的资源账本或评估协议反驳它。
\end{itemize}
\end{document}
